% GENERATED FILE. Edit canonical lessons, then run npm run generate:books.
% canonical-source-hash: fnv1a64:076dc0f832d7834c
% canonical-lessons: IT-C532-spazzare, IT-C532-spolverare, IT-C532-straccio, IT-C532-detersivo, IT-C532-aspirapolvere

\chapter{Le pulizie}
\label{ch:it-le-pulizie}
\hlchaptermodality{532}

\begin{chapteropening}
\textbf{By the end of this chapter:} \emph{I can talk about cleaning the house.}

Complete the last lesson of chapter 532: I can talk about cleaning the house.
\end{chapteropening}

\section[spazzare]{spazzare — to sweep}

\label{lesson:IT-C532-spazzare}



\begin{quote}
Before the new one: say the Italian for clothes peg, then the Italian for iron.
\end{quote}

\subsection*{You'll want to know: spazzare}
\textbf{spazzare} — \textquotedblleft{}to sweep\textquotedblright{}.

\begin{quote}
Dopo cena spazzo la cucina.
\end{quote}

\begin{quote}
\emph{After dinner I sweep the kitchen.}
\end{quote}

With \emph{la scopa} \textquotedblleft{}the broom\textquotedblright{}.

\begin{cousinweb}
\emph{spazzare} is usually traced to Late Latin \emph{spatiāre} (EtymDB), from \emph{spatium} \textquotedblleft{}space\textquotedblright{}: to clear a space.
\end{cousinweb}

\begin{grammarlens}[title={What you've built}]
One more word for telling, explaining and describing at B1.
\end{grammarlens}

\subsection*{Your turn}
\begin{itemize}
\raggedright
  \item \emph{Say it:} \emph{spazzare}
  \item \emph{Say it:} \emph{spazzare}, once more
  \item \emph{Say it:} a sentence of your own with \emph{spazzare}
  \item \emph{Recall:} say the Italian for clothes peg, then the Italian for iron, then say \emph{spazzare} again
\end{itemize}

\begin{tcolorbox}[breakable,colback=teal!4,colframe=teal!35!black,title={Before you move on}]
What does spazzare mean? (\textquotedblleft{}To sweep\textquotedblright{}.) Say it once more.
\end{tcolorbox}
\section[spolverare]{spolverare — to dust}

\label{lesson:IT-C532-spolverare}



\begin{quote}
Before the new one: say the Italian for iron, then the Italian for to sweep.
\end{quote}

\subsection*{You'll want to know: spolverare}
\textbf{spolverare} — \textquotedblleft{}to dust\textquotedblright{}.

\begin{quote}
Spolvero i mobili una volta alla settimana.
\end{quote}

\begin{quote}
\emph{I dust the furniture once a week.}
\end{quote}

Built on \emph{polvere} \textquotedblleft{}dust\textquotedblright{}.

\begin{cousinweb}
\emph{spolverare} is built from \emph{s-} and \emph{polvere} \textquotedblleft{}dust\textquotedblright{}, inherited from Latin \emph{pulvis, pulveris} (Lewis \& Short).
\end{cousinweb}

\begin{grammarlens}[title={What you've built}]
One more word for telling, explaining and describing at B1.
\end{grammarlens}

\subsection*{Your turn}
\begin{itemize}
\raggedright
  \item \emph{Say it:} \emph{spolverare}
  \item \emph{Say it:} \emph{spolverare}, once more
  \item \emph{Say it:} a sentence of your own with \emph{spolverare}
  \item \emph{Recall:} say the Italian for iron, then the Italian for to sweep, then say \emph{spolverare} again
\end{itemize}

\begin{tcolorbox}[breakable,colback=teal!4,colframe=teal!35!black,title={Before you move on}]
What does spolverare mean? (\textquotedblleft{}To dust\textquotedblright{}.) Say it once more.
\end{tcolorbox}
\section[lo straccio]{lo straccio — rag, cloth}

\label{lesson:IT-C532-straccio}



\begin{quote}
Before the new one: say the Italian for to sweep, then the Italian for to dust.
\end{quote}

\subsection*{You'll want to know: lo straccio}
\textbf{lo straccio} — \textquotedblleft{}rag, cloth\textquotedblright{}.

\begin{quote}
Pulisco il tavolo della cucina con uno straccio bagnato.
\end{quote}

\begin{quote}
\emph{I wipe the kitchen table with a wet cloth.}
\end{quote}

Old, torn clothes are \emph{stracci}.

\begin{cousinweb}
\emph{straccio} is built on \emph{stracciare} \textquotedblleft{}to tear\textquotedblright{}, usually traced to Vulgar Latin \emph{extractiāre}, from \emph{trahere} \textquotedblleft{}to drag\textquotedblright{} (Lewis \& Short).
\end{cousinweb}

\begin{grammarlens}[title={What you've built}]
One more word for telling, explaining and describing at B1.
\end{grammarlens}

\subsection*{Your turn}
\begin{itemize}
\raggedright
  \item \emph{Say it:} \emph{lo straccio}
  \item \emph{Say it:} \emph{lo straccio}, once more
  \item \emph{Say it:} a sentence of your own with \emph{lo straccio}
  \item \emph{Recall:} say the Italian for to sweep, then the Italian for to dust, then say \emph{lo straccio} again
\end{itemize}

\begin{tcolorbox}[breakable,colback=teal!4,colframe=teal!35!black,title={Before you move on}]
What does lo straccio mean? (\textquotedblleft{}Rag, cloth\textquotedblright{}.) Say it once more.
\end{tcolorbox}
\section[il detersivo]{il detersivo — detergent, washing-up liquid}

\label{lesson:IT-C532-detersivo}



\begin{quote}
Before the new one: say the Italian for to dust, then the Italian for rag.
\end{quote}

\subsection*{You'll want to know: il detersivo}
\textbf{il detersivo} — \textquotedblleft{}detergent, washing-up liquid\textquotedblright{}.

\begin{quote}
È finito il detersivo per i piatti.
\end{quote}

\begin{quote}
\emph{We have run out of washing-up liquid.}
\end{quote}

\emph{Il detersivo per la lavatrice} — washing powder.

\begin{cousinweb}
\emph{detersivo} is a learned word built on Latin \emph{dētersus}, from \emph{dētergēre} \textquotedblleft{}to wipe off\textquotedblright{} (Lewis \& Short).
\end{cousinweb}

\begin{grammarlens}[title={What you've built}]
One more word for telling, explaining and describing at B1.
\end{grammarlens}

\subsection*{Your turn}
\begin{itemize}
\raggedright
  \item \emph{Say it:} \emph{il detersivo}
  \item \emph{Say it:} \emph{il detersivo}, once more
  \item \emph{Say it:} a sentence of your own with \emph{il detersivo}
  \item \emph{Recall:} say the Italian for to dust, then the Italian for rag, then say \emph{il detersivo} again
\end{itemize}

\begin{tcolorbox}[breakable,colback=teal!4,colframe=teal!35!black,title={Before you move on}]
What does il detersivo mean? (\textquotedblleft{}Detergent, washing-up liquid\textquotedblright{}.) Say it once more.
\end{tcolorbox}
\section[l'aspirapolvere]{l'aspirapolvere — vacuum cleaner}

\label{lesson:IT-C532-aspirapolvere}



\begin{quote}
Before the new one: say the Italian for rag, then the Italian for detergent.
\end{quote}

\subsection*{You'll want to know: l'aspirapolvere}
\textbf{l'aspirapolvere} — \textquotedblleft{}vacuum cleaner\textquotedblright{}.

\begin{quote}
Passo l'aspirapolvere in tutte le stanze.
\end{quote}

\begin{quote}
\emph{I vacuum all the rooms.}
\end{quote}

\emph{Passare l'aspirapolvere} — to vacuum.

\begin{cousinweb}
\emph{aspirapolvere} is built from \emph{aspirare} \textquotedblleft{}to suck in\textquotedblright{}, from Latin \emph{aspīrāre} \textquotedblleft{}to breathe on\textquotedblright{}, and \emph{polvere} \textquotedblleft{}dust\textquotedblright{}, from Latin \emph{pulvis} (Lewis \& Short).
\end{cousinweb}

\begin{grammarlens}[title={What you've built}]
One more word for telling, explaining and describing at B1.
\end{grammarlens}

\subsection*{Your turn}
\begin{itemize}
\raggedright
  \item \emph{Say it:} \emph{l'aspirapolvere}
  \item \emph{Say it:} \emph{l'aspirapolvere}, once more
  \item \emph{Say it:} a sentence of your own with \emph{l'aspirapolvere}
  \item \emph{Recall:} say the Italian for rag, then the Italian for detergent, then say \emph{l'aspirapolvere} again
\end{itemize}

\begin{tcolorbox}[breakable,colback=teal!4,colframe=teal!35!black,title={Before you move on}]
What does l'aspirapolvere mean? (\textquotedblleft{}Vacuum cleaner\textquotedblright{}.) Say it once more.
\end{tcolorbox}
